%% kapa: exam=2 round=35 year=2024 period=3 subject=compensation_law subject_ko=감정평가 및 보상법규 seq=3 number=3 points=20
%% source: Q-Net 감정평가사 기출문제 (cst00309) artlSeq=5213532 file=raw/1008/5213532/2024년 제35회 감정평가사 2차시험 문제지.pdf
%% conversion: kordoc 4.24.51 -> tools/gen2.py (auto); review=proofread-sonnet; flags=none
\begin{longquestion}{【문제 3】}
A감정평가법인(이하 ‘A법인’이라 함)은 B민간임대아파트 분양전환대책위원회(이하 ‘B대책위원회’라 함)와의 용역계약에 따라 해당 아파트의 분양전환 가격산정을 위한 감정평가서를 제출하였다. B대책위원회는 임대사업자 X의 의뢰를 받은 Y감정평가법인의 감정평가 결과와 A법인의 감정평가 결과가 크게 차이가 나자 국토교통부장관에게 각 감정평가에 대한 타당성조사 실시를 요청하였고, 국토교통부장관은 한국감정원으로 하여금 타당성조사를 실시하도록 하였다. 한국감정원은 B임대아파트 분양전환 가격산정을 위한 감정평가가 모두 부적정하다는 타당성조사 결과를 국토교통부장관에게 통지하였다. 다음 물음에 답하시오. (20점) \par
\medskip\noindent 물음 1) 국토교통부장관은 타당성조사 결과에 근거하여 고의로 잘못된 평가를 한 A법인 소속 감정평가사 甲에 대하여 업무정지 6개월의 징계처분을 하였다. 이에 불복한 甲이 징계처분취소소송을 제기하였는바, 법원은 해당 징계처분을 업무정지 3개월의 징계처분으로 감경하는 판결을 할 수 있는지에 관하여 설명하시오. (10점) \par
\medskip\noindent 물음 2) 국토교통부장관은 고의로 잘못된 평가를 한 甲이 소속된 A법인에 대하여 성실의무에 위반하였다는 사유로 과징금부과처분을 하였다. A법인은 자신이 부담하여야 하는 성실의무를 충실히 이행하였다고 주장하며 과징금부과처분에 불복하고자 한다. 이때 A법인이 부담하는 성실의무의 내용을 설명하시오. (10점) \par
\end{longquestion}
